\subsection{Answers to the research questions}
\label{sec:research-answers}

The evidence now lets us return to the three questions in Section~\ref{sec:research-questions}. The answers connect the method, receiver validation, and quality comparisons without treating their different samples as one experiment.

\subsubsection{Question 1: How do the choices carry bits?}
The method assigns bits to indices in the ordered reply-location and comment-topic lists. A writer expresses those choices through an ordinary comment using the model's API, without token probabilities. Shared source texts, preparation and ordering rules, model settings, and the comment cutoff let the parties interpret the same indices; compressed messages also require the corresponding dictionary and representation. Multi-comment transport assembles a longer message across carriers under the agreed ordering. The shared codec checks and recorded example demonstrate the bit mapping and exact reconstruction, while Sections~\ref{sec:receiver-operation} and~\ref{sec:reconstruction-validation} explain receiver inputs and validation of list agreement.

\subsubsection{Question 2: How successful are generation and recovery?}
The implementation generates accepted carriers and recovers complete messages, with success depending on the payload and receiving conditions. Source-excerpt recovery was \OctExact/\OctAttempts{} under prepared shared context. The independent-payload development check recovered \DevExact/\DevAttempted{} attempted messages before stopping, with \DevUnattempted{} unattempted. With independently acquired archive context, the anchored seven-attempt policy recovered \FreshDevSevenExact/\FreshAssigned{} development and \FreshHeldSevenExact/\FreshAssigned{} untouched messages. These are separate populations. Together with the component checks, the traces distinguish transport operations from linguistic failures: topic confusion prevented recovery in complete generated sequences, while exhausted generation budgets left other messages incomplete. Sender acceptance and receiver extraction therefore remain separate outcomes. Sections~\ref{sec:october-recovery} and~\ref{sec:fresh-archive} give their conditions and denominators.

\subsubsection{Question 3: How do the comments compare with ZLG?}
The independent judge's answer depends on the criterion. On the selected 244-post set, our comments performed better on register-fit and were chosen less often as the weakest text, while ZLG~\cite{lin2024zero} was harder to identify as the inserted comment. Suspicion AUROC values were close. The comparison therefore supports specific quality findings rather than a single ranking across tasks. The selected sample and model judgments define the scope of those findings (Table~\ref{tab:llm-judge}). Separately, the assisted extraction described in Section~\ref{sec:baseline-choice} limits conclusions about independent message recovery. The quality tasks do not establish a general security advantage.
